\section{Lower-Prefix And Upper-Tail Survivor Families}
\label{sec:lower-prefix-upper-tail-branch}

The scale-barrier branch leaves a more refined obstruction than a generic high-frequency tail. When the dyadic decomposition is examined closely, two survivor families appear. One family lives below the active scale as lower-prefix memory. The other lives above the active scale as upper-tail leakage. The positive-forward aim of this branch is to collapse those families into one paid object or prove that one family pays the other.

The dyadic picture is useful. Fix an active scale \(j\). The velocity splits schematically as
\[
  u = u_{<j-C}+u_{\sim j}+u_{>j+C}.
\]
The lower-prefix part \(u_{<j-C}\) carries accumulated low-frequency strain and transport memory. The upper-tail part \(u_{>j+C}\) carries separated high-frequency leakage. The active shell \(u_{\sim j}\) feels both. A forward proof wants a single estimate of the form
\[
  \mathrm{LowerPrefix}_j(t)+\mathrm{UpperTail}_j(t)
  \leq \mathrm{paid\ quantity}_j(t),
\]
where the paid quantity is summable in the derivative norm needed for continuation.

The branch is promising because the two families seem complementary. Lower-prefix memory records what the active shell has inherited from larger-scale motion. Upper-tail leakage records what the active shell fails to contain above itself. If both are tied to the same cascade, one hopes for a conservation or monotonicity relation that prevents them from surviving independently.

The stopping point is separation. The lower-prefix family and the upper-tail family can be seen through the same terminal event, yet the estimates that control one do not automatically control the other. Lower-prefix control asks for active-window strain bounds. Upper-tail control asks for summable high-frequency transfer. A theorem that identifies them as one terminal object, or makes one pay the other, is an extra bridge.

A typical lower-prefix quantity has the schematic form
\[
  P_j^{\downarrow}(t)=\sum_{k<j-C} a_{j,k}\|P_k u(t)\|_{X_k},
\]
where the coefficients \(a_{j,k}\) measure how much low-scale content can influence the active shell. A typical upper-tail quantity has the form
\[
  T_j^{\uparrow}(t)=\sum_{k>j+C} b_{j,k}\|P_k u(t)\|_{Y_k}.
\]
The proof would close if one could show that both \(P_j^{\downarrow}\) and \(T_j^{\uparrow}\) are dominated by the same dissipative ledger. The branch stalls when the available ledger controls them in different orientations.

This gives the survivor of the branch: a split terminal family. The obstruction is no longer a single bad source term. It is a pair of possible terminal carriers whose relationship remains unsettled under positive-forward estimates. Lower-prefix memory can survive as past scale information. Upper-tail leakage can survive as future scale escape. The branch asks whether these are two views of the same object or two independent objects that require separate classification.

The CM-exit handoff is direct. A finite-time breakdown claim that uses either family has to select a terminal witness. If lower-prefix memory is the witness, it must still carry a positive packet, same-equation participation, and positive-scale field readout. If upper-tail leakage is the witness, it faces the same test. If the proof later identifies the two families as the same terminal object, the same test applies once to that object.

The branch conclusion for the rebuilt PDF is therefore provisional. The lower-prefix / upper-tail split is support for the CM-exit program until a bridge theorem unifies the families or assigns each family a first failed service. Its paper role is to show why branch organization is necessary: a terminal obstruction can have several faces before the membership test decides its final class status.